% BEGIN VOL07-EXPANSION VII30
\section{Supplementary worked examples}
\label{sec:vii30-pedagogy}

\begin{example}[Proper time on an inertial segment]\label{ex:vii30-ped-01}
In \(1+1\) Minkowski space, a timelike straight segment from \((0,0)\) to \((T,X)\) with \(T>|X|\) has proper time \(\sqrt{T^2-X^2}\).  This Lorentzian interval replaces Riemannian distance along timelike motion.
\end{example}

\begin{example}[Broken path has less proper time]\label{ex:vii30-ped-02}
For endpoints \((0,0)\) and \((2,0)\), the inertial path has proper time \(2\).  A symmetric broken path through \((1,a)\), \(|a|<1\), has total proper time \(2\sqrt{1-a^2}<2\), illustrating local maximality of timelike geodesics.
\end{example}

\begin{example}[Null line]\label{ex:vii30-ped-03}
The line \(x=t\) in \(1+1\) Minkowski space is a null geodesic.  Its proper time is zero even though it is nonconstant and geometrically meaningful, showing why null geometry cannot be reduced to a positive length minimization problem.
\end{example}

\section{Graded supplementary exercises}
The following exercises add a deliberate progression from concrete calculation to structural proof, hypothesis testing, application, and synthesis.

\subsection*{Standard computations and constructions}

\begin{exercise}[Interval type]\label{exr:vii30-ped-01}
Classify separation \((\Delta t,\Delta x)=(3,2)\).
\end{exercise}

\begin{hint}
Compute \(-9+4\).
\end{hint}

\begin{solution}
Timelike.
\end{solution}

\begin{exercise}[Null interval]\label{exr:vii30-ped-02}
Classify \((2,2)\).
\end{exercise}

\begin{hint}
Compute the interval.
\end{hint}

\begin{solution}
Null.
\end{solution}

\begin{exercise}[Spacelike interval]\label{exr:vii30-ped-03}
Classify \((1,2)\).
\end{exercise}

\begin{hint}
Compute \(-1+4\).
\end{hint}

\begin{solution}
Spacelike.
\end{solution}

\begin{exercise}[Proper time]\label{exr:vii30-ped-04}
Compute the proper time between \((0,0)\) and \((5,3)\) along the inertial segment.
\end{exercise}

\begin{hint}
Use \(\sqrt{T^2-X^2}\).
\end{hint}

\begin{solution}
It is \(4\).
\end{solution}

\begin{exercise}[Null proper time]\label{exr:vii30-ped-05}
What is the proper time along a null curve?
\end{exercise}

\begin{hint}
Its tangent norm-square is zero.
\end{hint}

\begin{solution}
Zero.
\end{solution}

\subsection*{Proofs}

\begin{exercise}[Causal type conservation]\label{exr:vii30-ped-06}
Prove the causal type of an affinely parametrized geodesic is constant.
\end{exercise}

\begin{hint}
Differentiate \(g(\dot\gamma,\dot\gamma)\).
\end{hint}

\begin{solution}
Metric compatibility and the geodesic equation make this quantity constant.
\end{solution}

\begin{exercise}[Time-orientability criterion]\label{exr:vii30-ped-07}
Explain why a nowhere-zero timelike vector field determines a time orientation.
\end{exercise}

\begin{hint}
Use it to choose one component of each timelike cone.
\end{hint}

\begin{solution}
Declare the cone containing the field vector to be future at each point; continuity gives a continuous choice.
\end{solution}

\begin{exercise}[Inertial maximality example]\label{exr:vii30-ped-08}
Show \(2\sqrt{1-a^2}<2\) for \(0<|a|<1\).
\end{exercise}

\begin{hint}
Square both sides.
\end{hint}

\begin{solution}
Since \(1-a^2<1\), the broken timelike path has smaller proper time than the inertial path.
\end{solution}

\begin{exercise}[Shared formulas]\label{exr:vii30-ped-09}
Why do Christoffel and curvature formulas have the same algebraic form in Riemannian and Lorentzian signatures?
\end{exercise}

\begin{hint}
Identify the property used in their derivation.
\end{hint}

\begin{solution}
They require a smooth nondegenerate metric and its inverse, not positive definiteness.
\end{solution}

\subsection*{Counterexamples and hypothesis tests}

\begin{exercise}[Geodesics minimize Lorentzian length]\label{exr:vii30-ped-10}
Test the direct Riemannian analogy for timelike curves.
\end{exercise}

\begin{hint}
Use the broken-path example.
\end{hint}

\begin{solution}
False.  Timelike geodesics locally maximize proper time under suitable conditions.
\end{solution}

\begin{exercise}[Null equals zero motion]\label{exr:vii30-ped-11}
Test the claim: zero proper time means a curve is constant.
\end{exercise}

\begin{hint}
Use a null line.
\end{hint}

\begin{solution}
False.  Nonconstant null curves have zero proper time.
\end{solution}

\begin{exercise}[Every manifold time-orientable]\label{exr:vii30-ped-12}
Test the claim for Lorentzian manifolds.
\end{exercise}

\begin{hint}
Time orientation is extra global structure.
\end{hint}

\begin{solution}
False.  Some Lorentzian manifolds admit no continuous global choice of timelike cone.
\end{solution}

\subsection*{Applications and investigations}

\begin{exercise}[Relativistic clock]\label{exr:vii30-ped-13}
Why does proper time model an ideal clock carried along a timelike worldline?
\end{exercise}

\begin{hint}
It is invariant under Lorentzian coordinate changes.
\end{hint}

\begin{solution}
The integral depends only on the spacetime metric and the worldline, giving an observer-independent elapsed time along that trajectory.
\end{solution}

\begin{exercise}[Causal reachability]\label{exr:vii30-ped-14}
What is the distinction between \(I^+(p)\) and \(J^+(p)\)?
\end{exercise}

\begin{hint}
Timelike versus causal curves.
\end{hint}

\begin{solution}
\(I^+(p)\) uses future timelike curves; \(J^+(p)\) allows both timelike and null curves.
\end{solution}

\subsection*{Challenge problems}

\begin{exercise}[Twin-path calculation]\label{exr:vii30-ped-15}
For endpoints \((0,0)\) and \((4,0)\), compute the proper time of the broken path through \((2,1)\).
\end{exercise}

\begin{hint}
Each leg has \(T=2\), \(|X|=1\).
\end{hint}

\begin{solution}
Each leg contributes \(\sqrt3\), so total proper time is \(2\sqrt3<4\).
\end{solution}

\begin{exercise}[Signature bridge]\label{exr:vii30-ped-16}
Identify the single algebraic change that turns a Riemannian metric into Lorentzian geometry locally.
\end{exercise}

\begin{hint}
Compare signatures.
\end{hint}

\begin{solution}
Positive definiteness is replaced by nondegeneracy with index one, introducing timelike, null, and spacelike directions.
\end{solution}

% END VOL07-EXPANSION VII30
